\section{Scientific Scope and Claim Boundary}
\label{sec:scope}

The study is an offline, public-data, development-only benchmark. It is not a
mission design authority, a flight certification exercise, or a performance
assessment of a quantum processor. A simulation result may be useful for
evaluating a model family, but it cannot by itself establish the performance
of a mission-owned trajectory system or a hardware quantum implementation.
The distinction is essential for a human-rated application, where human-system
requirements are separate from a surrogate's regression score
\cite{nasa3001}.

\subsection{Reader guide to labels and data stages}

Short identifiers are retained because they link the paper to immutable
repository records. They are audit keys, not scientific concepts that a reader
is expected to know in advance:
\begin{description}
\item[Gate.] A predefined authorization checkpoint. Gate 3 asks whether the
simulator is credible within its public-data scope; Gate 5 evaluates models on
development data; Gate 6 would be a separately authorized mission-level
experiment and was never opened.
\item[Decision and protocol IDs.] A label beginning with D identifies a
governance decision or bounded campaign, while P identifies a prospective
protocol. For example, D046 is the decision record for one successor test.
These labels are not data samples or performance grades.
\item[Model IDs.] C labels a classical candidate, Q labels a QML candidate,
and A labels an analysis control. Thus C06 is the physics-residual classical
control and Q01 is the primary quantum-kernel candidate.
\item[Figure IDs.] RFIG identifies the repository's provenance record for a
figure. It allows the plotted data and generating script to be found; it has no
statistical meaning.
\item[Development-only.] Development rows may be used to choose and compare
models. Calibration rows would adjust a frozen model's output, and final-test
rows would estimate performance only after all choices were fixed. This study
read neither calibration nor final-test rows for the reported campaigns.
\item[Fold, seed, and rung.] A fold is one group-preserving train/validation
partition. A seed is a fixed random-number starting point used to repeat a
model reproducibly. A rung is one sample-size stage in the prespecified
128--1,024-row screening sequence.
\end{description}

\subsection{Human-centered decision-support stance}

The proposed role of a model in this benchmark is limited to offline screening
information for qualified human review. It is not an autonomous controller,
an operational recommendation engine, or a replacement for mission analysis.
For an HCI-oriented reading, the central design requirements are traceability,
understandable scope, visible uncertainty, retained failures, and a clear
fallback to physical analysis. The evidence package supports those requirements
by showing the data boundary, simulator checks, controls, stopping decisions,
and excluded evidence pathways together with the metrics.

No human participants, operator interfaces, observational studies, or measures
of trust, workload, situation awareness, or decision quality were collected.
The paper therefore makes no empirical HCI performance claim. It instead
describes how a high-consequence AI research artifact can be structured so that
later human evaluation is not pre-empted by an unsupported model claim.

The evidence boundary was defined before reporting. Gate 3 supports simulator
credibility only for the eligible comparisons. Gate 5 supports a technical
decision for a frozen development benchmark. The subsequent projected-kernel
protocol (P001) and 16 bounded successor decisions (D034--D049) are
development-only follow-ups; they cannot retroactively change the Gate 5
decision. Calibration and final-test rows, mission-loop
scenarios, quantum-hardware or graphics-processor execution, threshold application to real
operations, trained-model release, and Gate 6 remain outside the study. Table
\ref{tab:claim-boundary} summarizes the resulting language constraints.

\begin{table}[t]
\caption{Claim boundary used throughout the manuscript. This table is derived from the committed claim-boundary record and does not add a new experimental result.}
\label{tab:claim-boundary}
\centering
\footnotesize
\begin{tabularx}{\textwidth}{@{}p{0.19\textwidth}X@{}}
\toprule
Evidence category & Permitted interpretation \\
\midrule
Gate 3: simulator credibility & Simulator agreement for eligible public-data checks, not flight validation. RTC3 lacked eligible post-event evidence, so it was not tested and is neither a pass nor a failure. \\
Gate 5: development benchmark & Under the frozen development benchmark, the tested QML candidates did not outperform strong classical controls. \\
Post-Gate follow-ups & Projected-kernel protocol P001 and successor decisions D034--D049 are development-only negatives. They do not revise Gate 5 or authorize Gate 6. \\
Screening lesson & The recall-first classical analysis is a future protocol-design lesson, not a deployed safety filter or safety result. \\
Prohibited claims & No quantum advantage, propellant saving, mission-loop validity, flight readiness, NASA approval, hardware behavior, final-test result, or QML invention claim. \\
\bottomrule
\end{tabularx}
\end{table}


Figure \ref{fig:timeline} records the governance state rather than a model
metric. It is included because a benchmark conclusion depends on which evidence
was authorized when. A later candidate may be scientifically informative, but
it does not become preregistered evidence merely because it is reported after a
negative primary result.

\begin{figure}[t]
\centering
\includegraphics[width=\textwidth]{figures/gate_decision_timeline.png}
\caption{Research governance through Gate 5 (RFIG-001). Gates 0--5 were
closed by human decisions; the technical Gate 5 result is a benchmark-specific
protocol FAIL, meaning its advancement hypothesis was rejected. Gate 6 remains
unopened. This is decision-state evidence, not a model
performance plot.}
\label{fig:timeline}
\end{figure}
